\section{Giới thiệu dữ liệu}

\subsection{Bối cảnh bài toán (Domain Situation)}
Netflix là một môi trường phân phối nội dung toàn cầu với danh mục phim đa quốc gia, đa ngôn ngữ và rất khác nhau về thể loại, mức độ phổ biến, điểm đánh giá, ngân sách và doanh thu. Khi số lượng phim lớn, người xem dashboard khó nắm được bức tranh tổng thể nếu chỉ đọc từng bản ghi riêng lẻ.

Nhóm chọn bài toán khám phá thư viện phim Netflix từ góc nhìn trực quan hóa dữ liệu. Đối tượng sử dụng dashboard là người học phân tích dữ liệu, người quản lý danh mục nội dung hoặc người cần khảo sát xu hướng phim toàn cầu. Họ muốn trả lời nhanh các câu hỏi như: nội dung tập trung ở quốc gia nào, ngôn ngữ nào chiếm ưu thế, chất lượng nội dung phân bố ra sao, mức đầu tư tài chính có đi kèm hiệu quả thương mại hay không, và cụm đạo diễn/diễn viên nào nổi bật.

\subsection{Nguồn gốc dữ liệu}
\begin{itemize}[leftmargin=2em]
    \item Tên bộ dữ liệu gốc: \texttt{netflix\_movies\_detailed\_up\_to\_2025.csv}.
    \item Vị trí trong repo: \texttt{data/raw/netflix\_movies\_detailed\_up\_to\_2025.csv}.
    \item Dữ liệu sau làm sạch dùng cho dashboard: \texttt{data/processed/movies.csv} cùng các bảng phụ \texttt{movie\_countries.csv}, \texttt{movie\_genres.csv}, \texttt{movie\_directors.csv}, \texttt{movie\_casts.csv}.
    \item Thời điểm nhóm sử dụng dữ liệu trong repo: 04/2026.
    \item Mục đích sử dụng: học thuật, phục vụ đồ án giữa kỳ môn Trực quan hóa dữ liệu.
\end{itemize}

\subsection{Mô tả tổng quan dữ liệu}
Dữ liệu chính ở dạng bảng, mỗi dòng biểu diễn một phim. Từ tệp gốc gồm 16{,}000 bản ghi và 18 thuộc tính, nhóm làm sạch còn 15{,}239 phim trong bảng xử lý cuối cùng. Bộ dữ liệu trải dài từ năm 2010 đến 2025, bao phủ 133 quốc gia và 68 ngôn ngữ được chuẩn hóa để phục vụ trực quan hóa đa chiều.

\begin{table}[H]
\centering
\renewcommand{\arraystretch}{1.3}
\begin{tabularx}{\textwidth}{|p{4cm}|X|}
\hline
\textbf{Thông tin} & \textbf{Nội dung} \\
\hline
Tên bộ dữ liệu & Netflix Movies Detailed up to 2025 \\
\hline
Nguồn trong repo & \texttt{data/raw/netflix\_movies\_detailed\_up\_to\_2025.csv} \\
\hline
Đơn vị quan sát & Mỗi dòng là một phim \\
\hline
Số lượng bản ghi & 16{,}000 (thô) / 15{,}239 (sau làm sạch) \\
\hline
Số lượng thuộc tính & 18 thuộc tính chính trong bảng \texttt{movies.csv} \\
\hline
Dạng dữ liệu & Dữ liệu bảng (tabular), kết hợp biến định danh, định tính, định lượng và mô tả văn bản \\
\hline
Miền thời gian & 2010--2025 \\
\hline
Mức phủ không gian & 133 quốc gia, 68 ngôn ngữ \\
\hline
\end{tabularx}
\caption{Bảng tóm tắt dữ liệu đầu vào}
\end{table}

\subsection{Tiền xử lý dữ liệu}
\begin{itemize}[leftmargin=2em]
    \item Chuyển các trường \texttt{release\_year}, \texttt{vote\_average}, \texttt{popularity}, \texttt{vote\_count}, \texttt{budget}, \texttt{revenue} về kiểu số để tính toán ổn định trong dashboard.
    \item Chuẩn hóa tên quốc gia và ngôn ngữ thông qua lớp dịch \texttt{TranslationService}, giúp hiển thị nhất quán giữa dữ liệu và giao diện tiếng Việt.
    \item Tách các cột danh sách như \texttt{genres}, \texttt{director}, \texttt{cast} thành các bảng bridge để phục vụ phân tích theo thể loại, đạo diễn và diễn viên.
    \item Giữ lại các biến tài chính cần thiết để tính ROI và so sánh ngân sách--doanh thu.
    \item Loại bỏ hoặc bỏ qua các trường không còn dùng trực tiếp trong dashboard, ví dụ \texttt{duration} và \texttt{language} ở tệp gốc.
\end{itemize}

\subsection{Bảng mô tả các biến quan trọng}
\begin{longtable}{|p{3cm}|p{2.8cm}|p{3cm}|p{5.2cm}|}
\hline
\textbf{Biến} & \textbf{Kiểu dữ liệu} & \textbf{Vai trò} & \textbf{Ý nghĩa} \\
\hline
\endfirsthead
\hline
\textbf{Biến} & \textbf{Kiểu dữ liệu} & \textbf{Vai trò} & \textbf{Ý nghĩa} \\
\hline
\endhead
\texttt{release\_year} & Ordered & Trục thời gian & Năm phát hành phim \\
\hline
\texttt{country} / bảng \texttt{movie\_countries} & Categorical & Bản đồ quốc gia & Quốc gia sản xuất hoặc phát hành được gán cho phim \\
\hline
\texttt{language\_name} & Categorical & Treemap ngôn ngữ & Ngôn ngữ hiển thị chính của phim \\
\hline
\texttt{genres} & Categorical (list) & Phân tích thể loại & Danh sách thể loại của phim; được tách thành bảng thể loại chi tiết \\
\hline
\texttt{vote\_average} & Quantitative & Chất lượng nội dung & Điểm đánh giá trung bình của phim \\
\hline
\texttt{popularity} & Quantitative & Mức độ quan tâm & Chỉ báo mức độ phổ biến của phim \\
\hline
\texttt{budget} & Quantitative & Phân tích tài chính & Ngân sách sản xuất \\
\hline
\texttt{revenue} & Quantitative & Phân tích tài chính & Doanh thu của phim \\
\hline
\texttt{director} / bảng \texttt{movie\_directors} & Categorical (list) & Góc nhìn sáng tạo & Đạo diễn của phim, dùng để tổng hợp doanh thu và điểm trung bình \\
\hline
\texttt{cast} / bảng \texttt{movie\_casts} & Categorical (list) & Góc nhìn sáng tạo & Dàn diễn viên, dùng để tìm nhóm kéo doanh thu \\
\hline
\end{longtable}
